\chapter{Tongue tie division infant}
\index{Tongue tie division infant}

\section*{Definition and Overview}
Ankyloglossia, colloquially known as tongue-tie, is a congenital craniofacial anomaly characterized by an abnormally short, thick, or tight lingual frenulum. The lingual frenulum is a midline fold of mucous membrane that connects the ventral surface of the tongue to the floor of the mouth and the alveolar ridge of the mandible. In infants, this structural variation can severely restrict the range of motion of the tongue, particularly its ability to protrude, elevate, and lateralize. This restriction often interferes with the complex, coordinated biomechanics of infant feeding, leading to significant difficulties in breastfeeding and, less frequently, bottle-feeding.

Tongue-tie division, medically termed frenotomy (or frenulotomy), is a minor, rapid surgical procedure performed to release the restricted frenulum. The primary objective of the procedure is to restore the physiological mobility of the tongue, thereby facilitating an effective, pain-free latch for the breastfeeding mother and ensuring adequate nutritional intake and growth for the infant.

Historically, the management of ankyloglossia has evolved through periods of shifting clinical paradigms. In previous centuries, midwives routinely divided the lingual frenulum of newborns using a sharpened fingernail or simple surgical tools immediately after birth to prevent potential speech and feeding issues. During the mid-twentieth century, with the rise of commercial infant formula and a decline in breastfeeding rates, interest in and diagnosis of tongue-tie decreased significantly, as bottle-feeding requires less active tongue mobilization than breastfeeding. However, the recent global resurgence in breastfeeding advocacy, supported by the World Health Organization (WHO) and the United Nations Children's Fund (UNICEF), has highlighted the critical role of normal tongue function in successful lactation. Consequently, the diagnosis and treatment of ankyloglossia have returned to the forefront of neonatal care, necessitating an interdisciplinary approach involving midwives, lactation consultants, pediatricians, pediatric dentists, and pediatric surgeons.

\section*{Epidemiology}
The epidemiological profile of ankyloglossia in neonates is characterized by varying prevalence rates and demographic patterns reported across clinical studies.
\begin{itemize}
  \item \textbf{Prevalence and Incidence:} The reported prevalence of ankyloglossia in infants ranges widely from 1\% to over 10\%. This substantial variation is primarily attributed to the historical lack of standardized, universally accepted diagnostic criteria and differences in the clinical specialties of the observers. In studies using validated assessment tools, the true prevalence of symptomatic ankyloglossia is estimated to be between 4\% and 8\% of all live births.
  \item \textbf{Gender Distribution:} There is a consistent male predominance in the incidence of ankyloglossia. Epidemiological data indicate that male infants are affected approximately two to three times more frequently than female infants. The reasons for this sex-linked discrepancy remain unknown, though it suggests sex-influenced genetic or epigenetic factors in the development of oral structures.
  \item \textbf{Familial and Genetic Patterns:} While the majority of ankyloglossia cases occur as isolated, sporadic anomalies, a significant familial predisposition is observed. Approximately 20\% to 25\% of affected infants have a positive family history of tongue-tie, suggesting hereditary transmission. Family pedigree studies have indicated that it can be inherited as an autosomal dominant trait with variable expressivity and penetrance.
  \item \textbf{Temporal and Regional Trends:} Over the past two decades, there has been a notable increase in the number of tongue-tie diagnoses and subsequent frenotomy procedures performed globally, particularly in high-income countries. This trend does not reflect a change in the biological incidence of the condition but is instead driven by increased clinical awareness, a greater emphasis on breastfeeding, and heightened parental education and concern regarding infant oral anatomy.
\end{itemize}

\section*{Aetiology and Pathophysiology}
The development of ankyloglossia is rooted in embryological processes occurring during the first trimester of gestation, leading to anatomical and functional alterations in the infant's oral cavity.
\begin{itemize}
  \item \textbf{Embryological Failure of Apoptosis:} During the fourth to seventh weeks of embryonic development, the tongue arises from the first, second, and third pharyngeal arches. In early development, the tongue is completely fused to the floor of the primitive oral cavity. As gestation progresses, a process of programmed cell death (apoptosis) and tissue remodeling occurs, creating the sublingual sulcus that separates the tongue from the floor of the mouth. The lingual frenulum remains as the only midline attachment. Ankyloglossia occurs when there is incomplete apoptosis or impaired tissue remodeling, leaving a thick, short, or anteriorly displaced fibrous attachment that fails to regress.
  \item \textbf{Genetic Mutations and Syndromes:} Although isolated ankyloglossia is typically non-syndromic, genetic mutations have been identified in specific familial forms. Mutations in the \textit{TBX22} gene, which encodes a T-box transcription factor, are associated with X-linked cleft palate and ankyloglossia. Additionally, ankyloglossia is a recognized clinical feature of several congenital syndromes, including Beckwith-Wiedemann syndrome, Pierre Robin sequence, Smith-Lemli-Opitz syndrome, and Orofaciodigital syndrome.
  \item \textbf{Impairment of Infant Suckling Biomechanics:} During normal breastfeeding, the infant must extend the tongue forward over the lower gum line, wrap the lateral borders of the tongue around the nipple-areola complex, and elevate the mid-tongue to the hard palate. This action creates a negative pressure vacuum that draws breast tissue into the mouth and initiates a peristaltic wave that pumps milk from the lactiferous ducts. When the lingual frenulum is restricted, the infant cannot protrude or elevate the tongue. This prevents the formation of an effective seal, leading to a loss of vacuum, audible clicking, and inefficient milk transfer.
  \item \textbf{Compensatory Mechanics and Maternal Trauma:} To compensate for the lack of tongue mobility, the infant relies on alternative, pathological muscle movements. The infant may clamp down with the lips and gums to secure the breast, resulting in friction and compression of the maternal nipple. This abnormal pressure prevents the nipple from reaching the safety of the soft palate, leading to severe nipple pain, skin shearing, cracks, bleeding, and localized vasospasm (Raynaud's phenomenon of the nipple).
  \item \textbf{Feedback Loop of Lactation Failure:} The combination of inefficient infant suckling and maternal pain establishes a negative feedback loop. Inadequate breast emptying due to poor milk transfer leads to breast engorgement, plugged ducts, and an increased risk of mastitis. Over time, the lack of sufficient mechanical stimulation of the nipple leads to a decrease in prolactin and oxytocin secretion, resulting in a secondary reduction in maternal milk supply and eventual premature weaning.
\end{itemize}

\section*{Clinical Features}
The clinical presentation of ankyloglossia encompasses anatomical anomalies of the tongue and functional consequences for both the infant and the breastfeeding mother.
\begin{itemize}
  \item \textbf{Anatomical Features of the Tongue:} The visual and physical examination of the infant's oral cavity may reveal several characteristic findings:
    \begin{itemize}
      \item A short, thick, fibrous, or inelastic lingual frenulum attaching near the tip of the tongue (anterior ankyloglossia).
      \item A posterior ankyloglossia, which is less visually obvious and presents as a submucosal, thick, or tight elastic band that restricts the base and mid-body of the tongue, detectable primarily through digital palpation.
      \item A notched, heart-shaped, or square-shaped appearance of the tongue tip during protrusion or elevation, caused by the tethering effect of the frenulum.
      \item Inability of the infant to protrude the tongue beyond the lower lip or lower gum line.
      \item Inability of the tongue to elevate toward the hard palate when the mouth is wide open (e.g., when the infant is crying).
    \end{itemize}
  \item \textbf{Infant Functional Symptoms:} The functional manifestations of restricted tongue movement during feeding include:
    \begin{itemize}
      \item Difficulty achieving or maintaining a deep, stable latch on the breast or bottle nipple, leading to frequent slipping off the breast.
      \item An audible clicking or smacking sound during feeds, indicating a repeated breakage of the oral vacuum.
      \item Prolonged feeding times (e.g., sessions lasting over 45 minutes) or conversely, very short, frustrated feeding attempts due to physical exhaustion.
      \item Poor weight gain, flat growth curves, or failure to thrive due to inadequate milk intake.
      \item Extreme irritability, fussiness, or falling asleep at the breast shortly after initiating a feed due to the high energy expenditure required to maintain a latch.
      \item Colic, excessive flatulence, and gastroesophageal reflux symptoms, secondary to aerophagia (swallowing excess air) caused by a poor oral seal.
    \end{itemize}
  \item \textbf{Maternal Symptoms:} The physiological impact on the breastfeeding mother is a critical component of the clinical picture:
    \begin{itemize}
      \item Severe, persistent nipple pain during and immediately after feeding, which does not resolve after the first week postpartum.
      \item Visible nipple trauma, including compression lines (a ``lipstick-shaped'' nipple after feeding), blanching, cracking, bleeding, bruising, or ulceration.
      \item Chronic breast engorgement, blocked milk ducts, and recurrent mastitis due to incomplete breast emptying.
      \item Psychological distress, anxiety regarding feeding sessions, and feelings of failure or guilt related to feeding difficulties.
    \end{itemize}
\end{itemize}

\section*{Differential Diagnosis}
When evaluating an infant with feeding difficulties and suspected ankyloglossia, it is essential to consider and rule out alternative or compounding diagnoses:
\begin{itemize}
  \item \textbf{Laryngomalacia:} This congenital softening of the supraglottic tissues is the most common cause of stridor in infants. It can cause significant coordination difficulties during swallowing and breathing, leading to feeding fatigue, coughing, sputtering, and poor weight gain, which can be mistaken for or coexist with tongue-tie.
  \item \textbf{Micrognathia and Retrognathia:} An abnormally small or posteriorly displaced lower jaw alters the spatial relationships within the oral cavity. This prevents the infant from positioning the lower gum line correctly under the areola, making a deep latch impossible regardless of tongue mobility.
  \item \textbf{Cleft Lip and Cleft Palate:} Congenital fissures of the lip, hard palate, or soft palate (including submucous cleft palate) prevent the formation of an airtight seal. This leads to nasal regurgitation, poor suction, and feeding failure.
  \item \textbf{Maternal Anatomical Variations and Pathology:} Feeding difficulties may arise from maternal factors alone, such as flat, inverted, or large nipples, which can make latching difficult for a newborn. Additionally, nipple infections like candidiasis (thrush) or bacterial colonization (e.g., \textit{Staphylococcus aureus}) can cause severe nipple pain and damage in the absence of infant anatomical anomalies.
  \item \textbf{Transient Neonatal Suction Dysfunction:} This refers to poor sucking coordination or weak suction in the absence of anatomical defects, often seen in late preterm infants, infants exposed to intrapartum maternal analgesia (e.g., epidural fentanyl or pethidine), or those recovering from birth trauma.
  \item \textbf{Upper Lip Tie (Maxillary Labial Frenulum):} A restricted upper lip frenulum can prevent the upper lip from flanging outward, causing a shallow latch and air swallowing. While frequently diagnosed alongside tongue-tie, its clinical significance and the necessity of its division remain subjects of debate.
\end{itemize}

\section*{When to Seek Medical Care}
Parents and caregivers should monitor their newborn's feeding patterns and physical development closely. It is appropriate to seek a non-urgent evaluation from a pediatrician, midwife, or certified lactation consultant if the mother experiences persistent nipple pain or damage, or if the infant struggles to latch, feeds continuously without satisfaction, or exhibits slow weight gain.

However, certain signs indicate severe feeding compromise, dehydration, or systemic illness, requiring immediate medical evaluation. Caregivers should contact emergency services, such as local emergency services in many countries, or proceed to the nearest emergency department if the infant:
\begin{itemize}
  \item Exhibits profound lethargy, is difficult to arouse for feeds, or does not respond to touch or sound.
  \item Shows signs of respiratory distress, including tachypnoea (rapid breathing), nasal flaring, grunting, intercostal or subcostal retractions, or cyanosis (bluish skin color, particularly around the lips and tongue).
  \item Demonstrates signs of severe dehydration, which include fewer than three wet diapers in a 24-hour period (for infants older than four days), a sunken fontanelle, dry or sticky mucous membranes, and a lack of tears when crying.
  \item Develops a fever with a rectal temperature of 38.0 degrees Celsius (100.4 degrees Fahrenheit) or higher, particularly in infants under three months of age.
  \item Experiences active, continuous, or excessive bleeding from the sublingual region following a frenotomy procedure that does not stop after applying pressure with a sterile gauze pad for several minutes.
\end{itemize}

\section*{Diagnosis and Investigations}
The diagnosis of ankyloglossia is fundamentally clinical, requiring a comprehensive structural and functional evaluation. Diagnostic testing, such as laboratory work or imaging, is rarely necessary.
\begin{itemize}
  \item \textbf{Clinical Breastfeeding Observation:} The cornerstone of diagnostic assessment is the direct observation of a feeding session by a trained healthcare provider. The clinician evaluates the infant's positioning, latch depth, sucking rhythm, the sound of milk transfer (listening for clicking or swallowing), and the mother's subjective pain levels.
  \item \textbf{Visual and Structural Inspection:} The clinician inspects the infant's oral cavity under direct light, observing the attachment points of the lingual frenulum on the ventral surface of the tongue and the alveolar ridge. The shape of the tongue during crying or attempts at protrusion is noted.
  \item \textbf{Digital Sublingual Palpation:} The clinician performs a manual examination by inserting a gloved finger into the infant's mouth, sweeping under the tongue. This allows the clinician to palpate the thickness, elasticity, and length of the frenulum, and to identify posterior tongue-ties that are not visible to the eye but can be felt as a tight, restricting band.
  \item \textbf{Hazelbaker Assessment Tool for Lingual Frenulum Function (HATLFF):} This is a validated, comprehensive scoring system used to determine the severity of tongue-tie and the clinical necessity of a frenotomy. It scores five anatomical items (e.g., length of frenulum, attachment point) and seven functional items (e.g., lift, protrusion, peristalsis). A function score of 10 or less, or an anatomy score of 8 or less when function is impaired, indicates that surgical division is warranted.
  \item \textbf{Bristol Tongue Assessment Tool (BTAT):} This is a simplified, reliable four-item assessment tool commonly used in busy clinical settings. It scores tongue appearance, attachment of the frenulum, tongue lift when crying, and tongue protrusion. A total score of 3 or less (out of 8) suggests a severe tongue-tie, while a score of 4 to 5 indicates moderate restriction that may require intervention if feeding is symptomatic.
\end{itemize}

\section*{Management}
The management of ankyloglossia in infants requires a multidisciplinary, stepped-care approach, prioritizing conservative measures before proceeding to surgical intervention.
\begin{itemize}
  \item \textbf{Conservative Lactation Support:} The initial step in management is intensive support from a lactation consultant or midwife. Interventions include optimizing breastfeeding positioning (such as the laid-back, cross-cradle, or football hold) to encourage a deeper latch, teaching latch-enhancement techniques, managing nipple pain with purified lanolin or hydrogel dressings, and using a breast pump to maintain maternal milk supply if the infant cannot transfer milk effectively.
  \item \textbf{Indication for Frenotomy:} Surgical division is indicated only when the infant has a diagnosed anatomical tongue-tie accompanied by persistent feeding difficulties (e.g., poor weight gain, severe maternal pain, or inability to latch) that have not resolved with conservative lactation support.
  \item \textbf{Frenotomy Procedure Technique:}
    \begin{itemize}
      \item \textbf{Preparation:} The procedure is typically performed in an outpatient clinic. The infant is securely swaddled to prevent movement, and an assistant stabilizes the head.
      \item \textbf{Anesthesia:} Due to the minimal nerve endings and vascularity of the lingual frenulum in early infancy, the procedure is often performed without local anesthesia. However, comfort measures are widely used, including the administration of oral sucrose solution (which triggers endorphin release) or the application of a small amount of topical anesthetic gel (such as 2\% lidocaine) to the sublingual area.
      \item \textbf{Surgical Cut:} The clinician uses a sterile metal director or gloved fingers to elevate the tongue, exposing the frenulum. Using sterile, blunt-ended surgical scissors, a single, quick cut is made through the midline of the frenulum close to the tongue, avoiding the submandibular ducts and deep veins. Alternatively, some practitioners use a soft-tissue laser (such as a CO2 or Erbium:YAG laser) to vaporize the tissue, which provides immediate cauterization.
      \item \textbf{Hemostasis:} Immediately after the cut, the infant is handed to the mother to breastfeed. Sucking and the natural composition of breast milk provide immediate pain relief, and the pressure of the tongue against the palate achieves rapid hemostasis. Any minor residual bleeding is managed with direct pressure using a sterile gauze pad.
    \end{itemize}
  \item \textbf{Post-Operative Exercises and Wound Care:} Caregivers are typically instructed to perform gentle active wound management to prevent premature healing and closure of the surgical site. This includes sublingual sweeps (running a clean, gloved finger along the diamond-shaped wound) and tongue stretching exercises (gently lifting the tongue toward the palate) 3 to 4 times daily for approximately two to three weeks.
\end{itemize}

\section*{Complications}
While infant frenotomy is a common and exceptionally safe procedure, it is associated with a small risk of complications that clinicians and parents must be aware of:
\begin{itemize}
  \item \textbf{Post-Procedural Bleeding (Hemorrhage):} Minor spotting or oozing is expected immediately after the procedure and during the first few feeds. Significant or prolonged bleeding is rare (occurring in less than 1\% of cases) and is typically managed with local pressure, topical adrenaline, or silver nitrate application.
  \item \textbf{Infection at the Surgical Site:} The oral cavity is highly vascularized and salivary proteins have antimicrobial properties, making infection extremely rare. However, if pathogenic bacteria colonize the wound, it can lead to localized swelling, erythema, purulent discharge, and systemic fever, requiring topical or systemic antibiotic therapy.
  \item \textbf{Premature Reattachment and Scarring:} The most common complication is the premature fusion of the cut edges of the frenulum, leading to recurrence of the tongue-tie. This occurs if post-operative stretches are not performed or if the wound heals with significant contracture, which may necessitate a repeat frenotomy.
  \item \textbf{Feeding Aversion and Oral Hypersensitivity:} A small number of infants may experience oral discomfort for several days following the procedure, leading to a temporary refusal to feed or a nursing strike. This is managed with supportive care, skin-to-skin contact, pain relief, and alternative feeding methods (such as syringe or cup feeding) to prevent dehydration.
  \item \textbf{Iatrogenic Injury to Sublingual Structures:} Inexperienced practitioners risk causing injury to surrounding structures, including the salivary glands (opening of Wharton's ducts) or the deep lingual veins, which can cause hematoma formation or salivary duct obstruction.
\end{itemize}

\section*{Prognosis and Outcomes}
The prognosis for infants undergoing frenotomy is outstanding. Clinical trials and cohort studies consistently report high rates of success, with approximately 80\% to 90\% of mothers reporting an immediate, subjective reduction in nipple pain and improved feeding comfort within 24 to 48 hours of the procedure.

Objective improvements in feeding mechanics are also observed, including a more stable latch, reduced aerophagia, and more efficient milk transfer, which correlates with improved infant weight gain trajectories. Long-term studies show that infants who undergo frenotomy when clinically indicated have significantly higher rates of breastfeeding continuation at three and six months postpartum compared to symptomatic infants who do not receive the procedure. The surgical wound, which forms a characteristic white or yellowish diamond-shaped patch under the tongue, represents normal granulation tissue and heals completely within 10 to 14 days without long-term structural or functional sequelae. The potential impact of early frenotomy on future speech development remains a subject of ongoing debate, but its primary, evidence-based indication in infancy remains the resolution of feeding difficulties.

\section*{Prevention and Self-Care}
Because ankyloglossia is a congenital, structural anomaly, there are no known prevention strategies to avoid its occurrence in utero. However, the secondary complications—such as lactation failure, maternal nipple injury, and infant malnutrition—can be effectively prevented and managed through early screening, intervention, and structured home care.
\begin{itemize}
  \item \textbf{Early Screening and Lactation Education:} Expectant parents should receive prenatal education regarding normal breastfeeding dynamics, latch indicators, and common challenges. Early postnatal screening of infant oral anatomy by midwives or pediatricians allows for timely identification of tongue-tie before severe feeding complications develop.
  \item \textbf{Wound Management and Stretches:} Following a frenotomy, parents must adhere to the prescribed post-operative wound care regimen:
    \begin{itemize}
      \item Perform sublingual stretches and sweeps as demonstrated by the clinician, ensuring the tongue is lifted to expose the diamond-shaped wound to prevent the edges from adhering.
      * Maintain a regular schedule of exercises, typically 3 to 4 times a day for two to three weeks, incorporating them into diaper changes or playtime to minimize infant distress.
    \end{itemize}
  \item \textbf{Post-Operative Pain Relief:} If the infant exhibits signs of discomfort (such as persistent crying, fussiness, or feeding refusal) in the first 24 to 48 hours post-procedure, parents may administer infant paracetamol (acetaminophen) under the guidance of a healthcare professional, using accurate weight-based dosing.
  \item \textbf{Nipple Recovery and Healing:} For the breastfeeding mother, self-care is vital for healing damaged nipples. Recommended practices include applying purified lanolin cream or breast milk to the nipples after feeds, using silver nursing cups to protect the skin from friction, wearing breathable cotton nursing bras, and seeking immediate treatment for any signs of bacterial or fungal nipple infection.
  \item \textbf{Frequent Feeding on Demand:} Encouraging frequent, short feeding sessions immediately after the procedure and in the subsequent days helps keep the tongue active and mobile, which aids in functional rehabilitation and prevents wound adhesion.
\end{itemize}

\section*{Sources and Further Reading}
\begin{itemize}
  \item \textbf{National Institute for Health and Care Excellence (NICE):} Division of ankyloglossia (tongue-tie) for breastfeeding (Interventional procedures guidance IPG149). London: NICE. Available at: \texttt{https://www.nice.org.uk/guidance/ipg149}
  \item \textbf{Academy of Breastfeeding Medicine (ABM):} Protocol \#11: Guidelines for the evaluation and management of neonatal ankyloglossia. ABM Clinical Protocol. Available at: \texttt{https://www.bfmed.org/protocols}
  \item \textbf{American Academy of Pediatrics (AAP):} Clinical report on the diagnosis and management of ankyloglossia in infants. Pediatrics Journal. Available at: \texttt{https://publications.aap.org/pediatrics}
  \item \textbf{Australian Breastfeeding Association (ABA):} Information guide on tongue-tie and breastfeeding. Melbourne: ABA. Available at: \texttt{https://www.breastfeeding.asn.au/bf-info/common-concerns-infant/tongue-tie}
  \item \textbf{NHS England:} Overview of tongue-tie, symptoms, and treatment options. London: National Health Service. Available at: \texttt{https://www.nhs.uk/conditions/tongue-tie/}
  \item \textbf{Mayo Clinic:} Patient care and health information on tongue-tie (ankyloglossia). Rochester: Mayo Foundation for Medical Education and Research. Available at: \texttt{https://www.mayoclinic.org/diseases-conditions/tongue-tie}
\end{itemize}